\section{Appendix N -- Proof Machinery}
% R005_APPENDIX_READER_FRAME
Reader frame. This appendix supports the main manuscript by explaining corpus role, evidence class, audit route, and limitation boundary. It is not a detached raw dump.


\label{sec:oc13-proof-machinery-appendix}

This appendix explains how proof obligations are read in the public monograph. It does not ask the reader to treat release-status rows as mathematical evidence. The proof route has four layers: a prose theorem statement with assumptions, a proof sheet with dependency references, a selected Lean or finite semantic witness where applicable, and a reopening condition that states what would make the wording too strong.


\subsection{Proof-route authority}

A promoted theorem-like statement is public only when the statement, assumptions, dependency route, and boundary are visible. If a later empirical replay or comparator result contradicts a claimed support route, the empirical or comparator repair is attempted first; if the contradiction persists, the theorem scope is narrowed or reopened. This order protects the model core without hiding failed evidence behind terminology.


\subsection{Minimality and irreducibility reading rule}

Minimality and K-level irreducibility are read as witness obligations. A component cannot be removed if removal changes the declared semantic verdict, loses a retained witness, breaks a liveness or boundary distinction, or collapses a theorem dependency. Conversely, an added observable may be lawfully demoted when it is inert under the declared verdict suite. The finite-model evidence package records executable cases for these distinctions.


\subsection{Comparator and compression boundary}

The release compares model-core structure against selected alternatives under declared assumptions. It does not claim that every possible alternative has been defeated. A comparator statement is valid only at the same claim class and cost boundary named in the comparator register. This keeps proof machinery separate from rhetorical unrestricted comparative claim.


\subsection{Reader checklist}

\begin{itemize}[leftmargin=1.8em]

\item \textbf{Theorem statement:} the public sentence must state the claim and its assumptions

\item \textbf{Dependency route:} the proof sheet, formalization reference, finite case, or replay row must be named where it carries support

\item \textbf{Boundary:} the counterexample or reopening condition must be explicit

\item \textbf{No status substitution:} a status word, release row, or package flag is never a proof

\end{itemize}
